\bookchapter{Kapitel 25: Gebrauchttelefone und digitales Erbe}

\chapterintro{Apples “Legacy-Kontakt”}

\booksection{: Vertrauenswürdige Kontakte benennen, die nach dem Tod Zugang zu iCloud-Daten beantragen können.}

Grundprinzip: Datenloeschung ist nicht „Loeschen“, sondern „Ueberschreiben”. Viele glauben, es reicht, Fotos in den Papierkorb zu verschieben und das Album zu leeren. Nach dem Weiterverkaufen scannen neue Besitzer mit Datenrettungs-Software – und alte Fotos, Chatverlaeufe und Ausweisfotos sind alle wieder da. Diese „Blamage“ passiert jedes Jahr.

\begin{enumerate}
  \item \textbf{: Frist festlegen, nach der Daten automatisch an bestimmte Kontakte gesendet oder das Konto gelöscht wird.}: Das letzte vollstaendige Backup erstellen, alle Konten abmelden (iCloud, Huawei-/Xiaomi-/Samsung-Konto, Google-Konto), „Mein Geraet suchen" deaktivieren – sonst wird der naechste Nutzer durch die Aktivierungssperre ausgesperrt.
  \item \textbf{: Die zuverlässigste Methode – wertvolle Daten regelmäßig auf eine lokale Festplatte herunterladen.}: Bei modernen Smartphones umfasst der „Werkseinstellungen zuruecksetzen" in der Regel ein verschluesseltes Loeschen – es zerstoert den Verschluesselungsschluessel selbst. Ist der Schluessel weg, werden alle Daten sofort zu unleserlichem Datenmuell.
  \item \textbf{Letzter „Taeuschungstrick" (nur bei alten Geraeten noetig)}: Moderne Smartphones (iOS 8+, Android 6+) sind standardmaessig vollverschlusselt (FBE); der Werksreset zerstoert den Verschluesselungsschluessel, sodass Daten theoretisch nicht wiederherstellbar sind.\textbf{Dieser Schritt ist nicht erforderlich}. Nur bei sehr alten, unverschluesselten Geraeten (vor 2015) ist es noetig, nach dem Werksreset den Speicher mit Muell-Dateien zu fuellen und dann erneut auf Werkseinstellungen zurueckzusetzen.
\end{enumerate}

\textbf{Schnellreferenz-Tabelle: Datenloesch-Pfade der Marken}: Obwohl das Prinzip aehnlich ist, versteckt jede Marke den „Werkseinstellungen zuruecksetzen“-Eintrag an etwas anderer Stelle. Die wichtigsten Unterschiede:

\begin{center}
\begin{tabularx}{\linewidth}{|X|X|X|}
\hline
\textbf{Marke} & \textbf{Loesch-Einstieg} & \textbf{Wichtige Unterschiede und Hinweise} \tabularnewline
\hline
Ultra & Einstellungen > System und Updates > Zuruecksetzen > Auf Werkseinstellungen zuruecksetzen & Zuvor Huawei-Konto abmelden und „Mein Geraet suchen" deaktivieren, sonst wird die Aktivierungssperre ausgeloest; im reinen Modus kann eine Passwortverifizierung noetig sein. \tabularnewline
Xiaomi / Redmi & Einstellungen > Mein Geraet > Auf Werkseinstellungen zuruecksetzen & Unten gibt es „Virtuelle SD-Karte formatieren" und „Alle Inhalte und Einstellungen loeschen" – beide ankreuzen; vor dem Abmelden des Xiaomi-Kontos „Geraet finden" deaktivieren. \tabularnewline
OPPO / OnePlus & Einstellungen > Weitere Einstellungen > Handy zuruecksetzen > Alle Daten vollstaendig loeschen & ColorOS bietet einen Modus „vollstaendig loeschen", der den Speicher mehrfach ueberschreibt; Bestaetigung per Sperrbildschirm-Passwort erforderlich. \tabularnewline
vivo & Einstellungen > Systemverwaltung > Backup und Reset > Alle Daten loeschen & Erst „Handyspeicher formatieren" ankreuzen, um den SD-Kartenbereich zu loeschen; vivo-Konto abmelden und „Handy finden" deaktivieren. \tabularnewline
Samsung & Einstellungen > Allgemeine Verwaltung > Zuruecksetzen > Auf Werkseinstellungen zuruecksetzen & Zuvor Samsung-Konto entfernen und „Mein Geraet suchen" deaktivieren; einige Modelle bieten die Option „automatisch ueberschreiben", die empfohlen wird. \tabularnewline
iPhone & Einstellungen > Allgemein > iPhone uebertragen oder zuruecksetzen > Alle Inhalte und Einstellungen loeschen & Apple ID muss zuvor abgemeldet werden, sonst wird der naechste Nutzer durch die iCloud-Aktivierungssperre blockiert; ab iOS 15 wird „iCloud-Daten beim Abmelden behalten" unterstuetzt. \tabularnewline
\hline
\end{tabularx}
\end{center}

\begin{tcbwarn}{Die drei am haeufigsten uebersehenen Dinge}
\begin{itemize}
  \item \textbf{Konto nicht abgemeldet}: Das Handy wird nach dem Verkauf fernggesperrt; man kommt selbst nicht mehr rein und muss schliesslich mit dem Ausweis zum offiziellen Kundendienst.
  \item \textbf{SD-Karte nicht entnommen}: Der Werksreset formatiert nicht die Fotos auf der externen SD-Karte – daran denken, sie manuell zu entnehmen oder separat zu formatieren.
  \item \textbf{USB-Stick/OTG-Geraet nicht getrennt}: Wenn beim Loeschen ein externes Geraet angeschlossen ist, koennen dessen Daten versehentlich geloescht werden oder der Loeschvorgang abbrechen.
\end{itemize}

\end{tcbwarn}

\booksection{Seriöse Recycling-Plattformen}

Die Daten auf einem Handy sind laengst nicht mehr „virtuell“. Sie beinhalten dein Geld, deine sozialen Beziehungen, deine Werke und sogar deine „digitale Identitaet”.\textbf{Wie sollen diese Vermoegenswerte nach dem Tod einer Person behandelt werden?}Das ist keine abstrakte Frage mehr, sondern eine echte Herausforderung fuer immer mehr Familien: Angehoerige kennen die Passwoerter nicht, Plattformen erlauben keine Vererbung, und Kryptowaehrungen sind fuer immer verloren.

\begin{tcbinfo}{Rechtliche Grundlage: Art. 49 des Gesetzes zum Schutz personenbezogener Daten}
Artikel 49 des 2021 in Kraft getretenen „Gesetzes der Volksrepublik China zum Schutz personenbezogener Daten“ legt fest:\textbf{»Stirbt eine natuerliche Person, koennen ihre naechsten Angehoerigen zur Wahrung ihrer eigenen rechtmassigen und berechtigten Interessen die in diesem Kapitel vorgesehenen Rechte auf Einsicht, Vervielfaeltigung, Berichtigung und Loeschung der personenbezogenen Daten des Verstorbenen ausueben; sofern der Verstorbene zu Lebzeiten anderslautend verfuegt hat.«}

Das bedeutet:

\begin{itemize}
  \item \textbf{Naechste Angehoerige haben eine Rechtsgrundlage}, um bei Plattformen die Bearbeitung der Konten und personenbezogenen Daten des Verstorbenen zu beantragen.
  \item „Rechtmassige und berechtigte Interessen" umfassen in der Regel Erbschaftsangelegenheiten, Kontoschliessung und Loeschung sensibler Privatsphaerendaten.
  \item Hat der Verstorbene zu Lebzeiten in einem Testament oder einer Vereinbarung anders verfuegt (z. B. einen Legacy Contact bestimmt), gilt diese Anordnung.
  \item Plattformen koennen nicht pauschal mit „Nutzerdatenschutz" ablehnen; die genauen Verfahren unterscheiden sich jedoch je Anbieter.
\end{itemize}

\end{tcbinfo}

\booksubsection{WeChat und Alipay: Konten sind nicht vererbbar, aber Guthaben schon}

Die Regeln der beiden in China am haeufigsten genutzten Apps sind unterschiedlich:

\begin{itemize}
  \item \textbf{WeChat-Konto}: Gemaess der „Tencent WeChat Software-Lizenz- und Servicevereinbarung" gehoert das Konto Tencent; der Nutzer hat nur ein Nutzungsrecht.\textbf{Das Konto selbst ist nicht vererbbar und nicht uebertragbar}. Aber das\textbf{Guthaben und Geld im Geldbeutel-/Finanzprodukt}innerhalbe des Kontos zaehlt zum persoenlichen Vermoegen. Naechste Angehoerige koennen mit Sterbeurkunde, Verwandtschaftsnachweis und Erbschein bei Tenpay Auszahlung oder Uebertragung beantragen.
  \item \textbf{Alipay}: Auch hier ist das Konto nicht direkt vererbbar, aber Alipay-Guthaben, Yuebao- und Finanzprodukte\textbf{zaehlen als vererbbares Vermoegen}. Erben koennen dies ueber die Alipay-Kundenservice-Hotline oder vor Ort mit Sterbeurkunde und Erbschein beantragen; Finanzprodukte der Ant-Anbieter erfordern den Weg ueber die Fondsgesellschaft.
  \item \textbf{Chatverlaeufe}: Rechtlich handelt es sich um Privatsphäre; die Plattform stellt sie den Angehoerigen nicht von sich aus zur Verfuegung. Bei Erbschaftshinweisen koennen sie ueber Klage oder Notar beantragt werden.
\end{itemize}

\begin{tcbwarn}{Das Smartphone verbindet sich direkt mit einem Kommunikationssatelliten in der Umlaufbahn und sendet Notrufnachrichten. Derzeit wird hauptsächlich das Senden von SMS und Standortinformationen unterstützt, nicht aber Telefonate.}
\begin{itemize}
  \item \textbf{Den Verstorbenen-Account nicht einfach mit Passwort an die Familie „weitergeben"}– das verletzt die Plattformvereinbarung. Wenn Guthaben eingefroren wird oder Risiko-Kontrollen ausgeloest werden, ist die Beweisfuehrung erschwert.
  \item \textbf{Nicht versuchen, mit dem Handy des Verstorbenen Verifizierungscodes abzurufen und „Guthaben abzuheben"}– das koennte als Diebstahl gewertet werden und rechtliche Konsequenzen nach sich ziehen.
  \item \textbf{Nicht zoegern}: Einige Konten werden nach langer Inaktivitaet automatisch geloescht; die Rueckgewinnung von Guthaben wird dann schwieriger.
\end{itemize}

\end{tcbwarn}

\booksubsection{Apple „Legacy Contact“: Fruehzeitig festlegen – doppelte Wirkung mit halbem Aufwand}

Ab iOS 15.2 bietet Apple die Funktion „Legacy Contact“ an. Nutzer koennen vertrauenswuerdige Kontakte vorab festlegen; nach dem Tod kann diese Person mit\textbf{Sterbeurkunde + Legacy-Contact-Zugangsschluessel}den Zugriff auf die iCloud-Daten des Verstorbenen beantragen (Fotos, Notizen, Kontakte usw.). Schritte:

\begin{enumerate}
  \item Am iPhone oeffnen:\textbf{[Einstellungen] > [Dein Name] > [Anmeldung und Sicherheit] > [Legacy Contact]}.
  \item Auf „Legacy Contact hinzufuegen" tippen, mit Face ID bestaetigen.
  \item Kontakt auswaehlen (aus dem Adressbuch); das System generiert einen\textbf{Zugangsschluessel}.
  \item Den Zugangsschluessel per SMS, E-Mail oder als Ausdruck uebergeben und um sichere Aufbewahrung bitten.
  \item Mehrere Kontakte als Backup hinzufuegen, um Single-Point-of-Failure zu vermeiden.
\end{enumerate}

Der Eingetragene benoetigt kein Apple-Geraet; er kann mit Schluessel und Sterbeurkunde auf der Apple-Website den Zugriff beantragen. Der Vorgang dauert in der Regel 7 Tage.

\booksubsection{Google „Inaktivitaetskonten-Manager“: Automatische Regelung des digitalen Nachlasses}

Google bietet den „Inaktivitaetskonten-Manager (Inactive Account Manager)“ an, der bei lang andauernder Inaktivitaet vordefinierte Aktionen ausfuehrt – eines der ausgereiftesten Mechanismen fuer digitale Nachlassverwaltung.

\begin{enumerate}
  \item Repository-Namen (Repository name) eingeben, z. B.\textbf{myaccount.google.com/inactive}, mit dem Google-Konto anmelden.
  \item Timeout-Zeitraum festlegen: 3 / 6 / 12 / 18 Monate ohne Anmeldung bis zur Ausloesung.
  \item Vertrauenswuerdige Kontakte hinzufuegen (max. 10) und festlegen, welche Daten jeder zugreifen darf (Gmail, Fotos, Drive usw.).
  \item Optional: Automatische Antwort-E-Mail einrichten, die Absender informiert, dass du nicht mehr antworten kannst.
  \item Letzte Auswahl: Nach Ablauf entweder\textbf{Daten mit Kontakten teilen}oder\textbf{Konto direkt loeschen}.
  \item Handynummer als Absicherung eintragen, um Fehlausloesungen zu vermeiden.
\end{enumerate}

\booksubsection{Sozialnetzwerke in China: Gedenkmodus und Loeschung}

Die grossen chinesischen Social-Media-Plattformen gehen auf zwei Arten mit verstorbenen Nutzern um:\textbf{„Gedenkmodus“ (Konto fuer Angehoerige erhalten)}„\textbf{„Loeschung” (komplett entfernen)}. Angehoerige koennen mit Sterbeurkunde, Verwandtschaftsnachweis und eigenem Ausweis beim Kundenservice der Plattform beantragen.

\begin{itemize}
  \item \textbf{Weibo}: Unterstuetzt den „Gedenkmodus"; Angehoerige stellen einen Antrag mit Sterbeurkunde. Das Konto wird gesperrt, der Inhalt bleibt erhalten, und das Profil zeigt „Dieses Konto wurde als Gedenkkonto eingerichtet." Eine Loeschung ist mit Erbschein moeglich.
  \item \textbf{Douyin (TikTok China)}: Ueber „Ich – Einstellungen – Feedback und Hilfe – Konto und Sicherheit – Konto loeschen" oder manuellen Kundenservice; Sterbeurkunde, Verwandtschaftsnachweis und Ausweis des Antragstellers einreichen. Nach Pruefung kann das Konto geloescht oder in den Gedenkmodus versetzt werden.
  \item \textbf{Bilibili}: Ueber den Kundenservice Sterbeurkunde und Verwandtschaftsnachweis einreichen; das Konto kann geloescht oder als Gedenkseite erhalten werden. Die Videos eines UP-Hosts werden unter Beruecksichtigung von Urheberrechtsangaben behandelt.
  \item \textbf{Xiaohongshu}: Ueber den offiziellen Kundenservice; mit Sterbeurkunde und Verwandtschaftsnachweis kann die Loeschung des Kontos und der Inhalte beantragt werden.
\end{itemize}

\begin{tcbinfo}{Wichtige Hinweise}
\begin{itemize}
  \item Plattformregeln werden aktualisiert; vor der Beantragung die aktuelle Anleitung des Kundenservice pruefen.
  \item Alle Nachweisdokumente zuerst einscannen und sichern, bevor sie eingereicht werden; Originale sorgfaeltig aufbewahren.
  \item Bei mehreren Plattformen kann ein einheitlicher\textbf{Erbschein}beantragt werden, um wiederholte Gaenge zu ersparen.
\end{itemize}

\end{tcbinfo}

\booksubsection{Die Schwierigkeit der Vererbung digitaler Vermoegenswerte}

Neben Social-Media-Konten besitzen immer mehr Menschen\textbf{wirklich wertvolle digitale Vermoegenswerte}, deren Vererbung komplizierter ist:

\begin{itemize}
  \item \textbf{Kryptowaehrungen (Bitcoin, Ethereum usw.)}: Nach dem Prinzip „privater Schluessel = Eigentum" gilt:\textbf{Ohne privaten Schluessel sind die Vermoegenswerte verloren}. Erben koennen nicht „per Klage" auf der Blockchain intervenieren. Empfehlung: Recovery Phrase / private Schluessel aufgeteilt an mehreren physisch sicheren Orten aufbewahren (z. B. Bankfach + Haustresor) und den Zugriffsweg im Testament beschreiben.
  \item \textbf{Boersen-Konten}: Aehnlich wie bei Bankkonten koennen Erben mit Erbschein bei der Boerse (z. B. Binance, OKX) einreichen; grenzueberschreitende Boersen haben jedoch lange Bearbeitungszeiten.
  \item \textbf{Domains}: Zaehlen als vererbbares Vermoegen; Erben koennen nach Einreichung des Erbscheins beim Registrar (z. B. Alibaba Cloud, GoDaddy) die Uebertragung beantragen.
  \item \textbf{Geistiges Eigentum (Originalartikel, Videos, Patente, Software-Urheberrechte)}: Die Vermoegensrechte aus dem Urheberrecht sind vererbbar (Persoenlichkeitsrechte nicht). Erben koennen die Ertragrechte bis 50 Jahre nach dem Tod des Urhebers wahrnehmen.
  \item \textbf{Cloud-Speicher und Online-Festplatten}: Inhaltsdaten gehoeren der Person und sind vererbbar; einige Plattformvereinbarungen sehen jedoch vor, dass das Konto selbst nicht uebertragbar ist – erst herunterladen, dann loeschen.
\end{itemize}

\textbf{Der wahre Preis „verschwindender“ Kryptowaehrungen}\par

Laut einem Bericht von Chainalysis sind weltweit bereits\textbf{ueber 3 Millionen Bitcoin}durch verlorene private Schluessel unwiederbringlich verloren – das entspricht einem Sechstel der vorhandenen Gesamtmenge. Wenn du Krypto-Vermoegen haeltst, stelle sicher, dass mindestens eine Person, der du absolut vertraust, den\textbf{Wiederherstellungsweg}kennt – sonst schlummern diese Vermoegenswerte fuer immer auf der Blockchain.

\booksubsection{Eine Vorlage fuer ein „digitales Testament“}

Ein herkoemmliches Testament deckt digitale Vermoegenswerte nicht ab. Es wird empfohlen, jaehrlich einen Anhang „Digitales Testament“ zu aktualisieren und zusammen mit dem offiziellen Testament aufzubewahren. Die folgende Vorlage dient als Referenz:

\textbf{„Anlage zur Regelung persoenlicher digitaler Vermoegenswerte“ (Vorlage)}\par

\textbf{Teil 1: Kontoliste}(Nur Kontoname und Plattform erfassen;\textbf{keine Passwoerter notieren})

\begin{itemize}
  \item Soziale Netzwerke: WeChat (Konto xxx), Weibo, Douyin, Bilibili, Xiaohongshu usw.
  \item Zahlung: Alipay, WeChat Pay, einzelne Banking-Apps
  \item Cloud: iCloud, Huawei Cloud, Xiaomi Cloud, Google-Konto, Microsoft-Konto
  \item E-Mail: QQ-Mail, NetEase-Mail, Gmail usw.
  \item Sonstige: Online-Festplatten, Abodienste, Spielekonten usw.
\end{itemize}

\textbf{Teil 2: Passwort-Verwahrung}

\begin{itemize}
  \item Master-Passwort: In einem Passwort-Manager (wie Bitwarden, 1Password) verwahren; das Master-Passwort auf Papier notieren und im Haustresor verschliessen.
  \item Wiederherstellungsschluessel / Recovery Phrase: Aufgeteilt an mindestens zwei physischen Orten aufbewahren (z. B. Bankfach + Haustresor).
  \item Berechtigte Person: Name, Ausweisnummer, Verwandtschaftsverhaeltnis, Kontaktdaten.
\end{itemize}

\textbf{Teil 3: Verfuegungswillen}

\begin{itemize}
  \item Sollen die jeweiligen Konten in den „Gedenkmodus" versetzt oder „geloescht" werden?
  \item Sollen fuer iCloud-/Google-Daten ein „Legacy Contact" oder der „Inaktivitaetskonten-Manager" eingerichtet werden?
  \item Wer erbt die Krypto-Vermoegenswerte? Zugangsweg zum privaten Schluessel beschreiben.
  \item Wer erhaelt die Ertraege aus urheberrechtlich geschuetzten Werken?
\end{itemize}

\textbf{Teil 4: Kontaktpersonen und Plattformen}

\begin{itemize}
  \item Vertrauenswuerdiger Testamentsvollstrecker: Name, Kontaktdaten, Ersatzkontakt.
  \item Kundenservice-Zugaenge der relevanten Plattformen und Liste der erforderlichen Nachweisdokumente.
  \item Aktualisierungsdatum dieser Anlage: \_\_\_\_/\_\_/\_\_; naechste Ueberpruefung: \_\_\_\_/\_\_/\_\_.
\end{itemize}

\begin{tcbwarn}{Rechtlicher Hinweis zu „Passwoertern"}
Passwoerter direkt im Testament zu notieren birgt das Risiko der Preisgabe; nach der Testamenteroeffnung koennen alle Erben sie einsehen.\textbf{Sicherer ist: Im Testament nur „Kontoliste + Verfuegungswillen + Aufbewahrungsort der Passwoerter“ notieren; die eigentlichen Passwoerter ueber das Master-Passwort des Passwort-Managers separat an den Vertrauensvollstrecker uebergeben}. So wird dem Recht genuegt und das Risiko der Passwortpreisgabe gesenkt.

\end{tcbwarn}

\booksection{(Optional) Den Speicher mit Mülldateien füllen und noch einmal zurücksetzen.}

Ein Smartphone enthaelt Dutzende Metalle wie Gold, Silber, Kupfer und Seltene Erden. Unueberlegt weggeworfen verschmutzt es den Boden und verschwendet Ressourcen. Es gibt drei richtige Recycling-Kanaele mit unterschiedlichen Schwerpunkten:

\begin{center}
\begin{tabularx}{\linewidth}{|X|X|X|X|}
\hline
\textbf{Kanal} & \textbf{Ruecknahmewert} & \textbf{Datensicherheit} & \textbf{sRGB} \tabularnewline
\hline
Offizielle Inzahlungnahme (Huawei/Apple/Xiaomi usw.) & Mittel, oft ergaenzt durch Gutscheine & Am hoechsten; offizieller, standardisierter Ablauf & Am guenstigsten, wenn ohnehin ein Geraet derselben Marke gekauft wird \tabularnewline
Drittanbieter-Recycling-Plattform (Aihuishou, Zhuanzhuan, Huisoubao usw.) & Etwas hoeher; Preise vergleichbar & Hoch; mit Datenloesch-Bericht & Wer direkt Geld haben will, ohne Markenbindung \tabularnewline
Spende (z. B. an Laendschulen, Wohlfahrtsorganisationen) & Keine finanzielle Rueckzahlung & Selbstaendige vollstaendige Datenloeschung erforderlich & Geraet noch funktionstuechtig; Gutes tun moechten \tabularnewline
Weiterverwendung (selbst behalten und umwidmen) & Spart Geld fuer ein Neugeraet & Der urspruengliche Besitzer nutzt weiter; kein Risiko & Geraet noch leistungsfaehig; Bedarf fuer Umwidmung im Haushalt \tabularnewline
\hline
\end{tabularx}
\end{center}

\begin{itemize}
  \item \textbf{Die nächste Station wählen: offizielles Recycling, seriöse Plattform oder Verschenken.}: Die empfehlenswerteste Methode: standardisierter Ablauf mit garantierter Datensicherheit. Apple, Huawei, Xiaomi und andere Marken haben auf ihrer Website einen Eingang „Inzahlungnahme"; nach Bewertung wird der Betrag vom Neupreis abgezogen.
  \item \textbf{Die "Auf-Null-setz"-Zeremonie beim Abschied vom alten Smartphone und die Planung des "Nachlasses" der digitalen Bestände sind nicht nur ein würdiges Ende, sondern eher ein Brief an die Zukunft. Lass uns heute ein wenig die Fantasie schweifen lassen und auf die am Horizont aufziehenden Stürme blicken.}: Waehle eine renommierte grosse Recycling-Plattform. Nach der Bestellung holt ein Kurier das Geraet ab; nach der Pruefung wird der Betrag ueberwiesen. Waehle unbedingt eine Plattform, die\textbf{einen „Datenloesch-Bericht" anbietet}, und bewahre die Bestaetigung auf.
  \item \textbf{Spende}: Kontaktiere oertliche Laendschulen oder Gemeinnuetzige Organisationen; alte Handys koennen als Unterrichts-Tablet oder Seniorenhandy umgebaut werden.\textbf{Vor der Spende unbedingt den vollstaendigen Loeschvorgang durchfuehren und alle SIM-/SD-Karten entfernen}.
  \item \textbf{Faltbare Smartphones}: Zum Auto-Navigator, Smart-Home-Zentrale, Ueberwachungskamera, Babyphone oder digitalen Bilderrahmen umbauen. Ein altes Handy kann noch Jahre dienen.
\end{itemize}

\begin{tcbinfo}{Die Wahrheit ueber Elektroschrott}
Die Leiterplatte eines Handys enthaelt\textbf{Gold, Silber, Palladium, Kupfer}und weitere Edelmetalle; jede Tonne alter Leiterplatten enthaelt ein Vielfaches des Goldgehalts von Golderz. Werden sie jedoch achtlos weggeworfen, gelangen Blei, Quecksilber und Cadmium in Boden und Grundwasser und vergiften diese jahrzehntelang.\textbf{Richtiges Recycling schuetzt die Umwelt und laesst seltene Metalle im Kreislauf zirkulieren.}

\end{tcbinfo}

\begin{tcbpractice}{AR-Brillen}
\begin{enumerate}
  \item Letztes vollstaendiges Backup: Fotos in die Cloud hochladen, wichtige Dateien auf den PC kopieren, Chatverlaeufe exportieren und archivieren.
  \item Alle Konten abmelden: iCloud / Huawei / Xiaomi / Google; „Mein Geraet suchen" deaktivieren.
  \item SIM- und SD-Karte entnehmen, alle OTG-Geraete trennen.
  \item Verschluesselten Werksreset durchfuehren; gegebenenfalls „Speicher formatieren" manuell ankreuzen.
  \item (Nur alte Geraete) Speicher mit Muell-Dateien fuellen, dann erneut auf Werkseinstellungen zuruecksetzen.
  \item Die naechste Station waehlen: offizielle Inzahlungnahme, seriocese Recycling-Plattform, Spende oder Verschenken an die Familie.
  \item Die „Liste der digitalen Vermoegenswerte" aktualisieren und ausgeschiedene Geraete streichen.
\end{enumerate}

\end{tcbpractice}

\begin{tcbpractice}{Zielgerichtete Speicherbereinigung}
\begin{enumerate}
  \item Fuer dein aktuelles Hauptgeraet einen Apple „Legacy Contact" einrichten oder den Google „Inaktivitaetskonten-Manager" aktivieren.
  \item Mit der Vorlage in diesem Abschnitt eine „Anlage zur Regelung persoenlicher digitaler Vermoegenswerte" entwerfen und den Aufbewahrungsort mit der Familie besprechen.
  \item Pruefen, ob zu Hause 2 oder mehr ungenutzte alte Handys liegen; eine Inzahlungnahme oder Recycling bei einer offiziellen/seriocese Plattform vereinbaren.
\end{enumerate}

\end{tcbpractice}
